% section: 特殊関数型

\section{特殊関数型}
ここまで,様々な漸化式について学んできた.
しかしながらこれまでに説明した方法ではで解けない漸化式も多々ある.
例えば対数型から1例あげよう.
\[
  a_{n+1}=a_n^2
\]
これは解ける.
しかし
\[
  a_{n+1}=a_n^2+1
\]
これは解けない.なぜならば対数をとってしまうと$+1$の項が処理できないからである.
だが,
\[
  a_{n+1}=a_n^2+2
\]
は解ける.
\[
  a_n=\alpha^{n}+\alpha^{n-1}
\]
と置くことで,帰納法を用いて示すことができる.
これが数学の大変興味深く面白い部分である.
対数型では積の形を和の形にできるという,対数の強みを使って解いた.
その結果,和が含まれる形には極端に弱い解法になってしまったわけだ.

この章では,一般には解けない形状をした式だが,ある特定の,数学的に特殊な性質を持っている故にとけるものを扱っている.
いわば運良く偶然解ける漸化式である.
\subsection{三角関数型}
この漸化式について考えよう.
\[
  a_1=\frac{\sqrt{3}}{2},\quad a_{n+1}=2a_n^2-1
\]
勘のいい,きちんと学んできた人には$\cos$の倍角公式が想起されるだろう.
実際初項も思わせぶりな値である.
念のため,三角関数の復讐をしておこう.
一般に以下のことが成り立つ.
\begin{align*}
  \sin (\alpha+\beta) & =\sin \alpha \cos \beta +\cos \alpha + \sin \beta                                                      \\
  \cos (\alpha+\beta) & =\cos \alpha \cos \beta -\sin \alpha + \sin \beta                                                      \\
  \tan (\alpha+\beta) & =\frac{\sin (\alpha+\beta)}{\cos (\alpha+\beta)}=\frac{1+\tan\alpha+\tan \beta}{1-\tan\alpha\tan\beta}
\end{align*}
また,\,$\alpha=\beta$の時,
\begin{align*}
  \sin 2x =2\sin x \cos x \\
  \cos 2x =2\cos^2 -1     \\
  \tan 2x =\frac{1+2\tan x}{1+\tan^2 x}
\end{align*}
これらを倍角公式と呼ぶ.また,このことから
\begin{align*}
  \sin^2 x= \frac{1-\cos x}{2} \\
  \cos^2 x= \frac{1+\cos x}{2}
  \tan^2 x=\frac{1-\cos x}{1+\cos x}
\end{align*}
加えて,ここでは更に3倍角公式についても扱う.
\begin{align*}
  \sin 3x = 3\sin x - 4 \sin^3 x
  \cos 3x = -3\cos x + 4\cos^3 x
\end{align*}


\subsection{基本対称式型}
この章の初めに,
\[
  a_{n+1}=a_n^2+2
\]
は解けるということを紹介した.
同じ理屈で
\[
  a_{n+1}=a_n^3+3a_n
\]
も解ける.
その手法について説明していこう.
\subsection{ニュートン法型}
漸化式の解説に入る前に,一次近似について復習しておこう.
\[
  a_{n+1}=\frac{a_n^2+2}{2a_n}
\]
\subsection{初項依存型}
まずはこの漸化式を見てもらいたい.
\[
  a_{n+1}=\frac{a_n^2-1}{2},a_1=c=2
\]
